\chapter{From a hand-made formula to a learned one}
\label{ch:learned}

\begin{objectives}
\begin{itemize}
\item Turn the hand-tuned evaluation $f(s) = \sum_i w_i\phi_i(s)$ of Chapter~\ref{ch:problem}
  into a formula whose weights are found from data.
\item Understand a derivative as a slope, and perform two steps of gradient descent by hand.
\item Say why a probability output needs the logistic function and a different loss.
\item \emph{Prove} that a linear formula cannot express "opposite-side castling is sharp",
  and then build the two-neuron network that can.
\item Perform a forward pass through a small network by hand and get the right answer.
\end{itemize}
\end{objectives}

\section{The only change that matters}

Chapter~\ref{ch:problem} left us with
\[
  f(s) \;=\; \sum_{i=1}^{k} w_i\,\phi_i(s) ,
\]
features $\phi_i$ chosen by a person, weights $w_i$ tuned by a person. Part~\textsc{iv} is
about replacing "by a person" with "by data", twice: first the weights (this chapter), then
the features themselves (Chapter~\ref{ch:attention}).

Nothing exotic arrives. A neural network is this same equation, applied several times in a
row, with a non-linear squash in between. If you understand this chapter, the only genuinely
new idea left in the whole of Part~\textsc{iv} is \emph{attention}, and that is one chapter.

\section{Learning the weights}

\subsection{The data and the target}

Take a large pile of positions in which the eventual result is known --- for Leela, positions
from its own self-play games. Each example is a pair: a position $s$ and a target $y$, the
result from the side to move's point of view ($+1$ win, $0$ draw, $-1$ loss).

We want weights that make $f(s)$ close to $y$ across the pile. "Close" needs a definition;
the standard one is squared error, summed over $M$ examples:
\begin{equation}
  \label{eq:mse}
  L(w) \;=\; \frac{1}{M}\sum_{m=1}^{M}\bigl(f(s_m) - y_m\bigr)^2 .
\end{equation}
Squaring does two jobs: it makes errors positive regardless of direction, and it punishes
large errors disproportionately (being wrong by $0.8$ is sixteen times worse than being wrong
by $0.2$, not four times).

\begin{notationbox}
$L$ is the \textbf{loss}: one number saying how badly the current weights fit the data. All
of training is the search for weights that make it small. "Training a network" and
"minimising $L$" are the same sentence.
\end{notationbox}

\subsection{Which way is downhill?}

We need to change $w$ so $L$ falls. The tool is the \textbf{derivative}, and for our purposes
it needs no more definition than this:

\begin{keyidea}
$\dfrac{\partial L}{\partial w_i}$ is the \textbf{slope} of the loss with respect to weight
$i$: how much $L$ changes if you nudge $w_i$ up by a tiny amount. If the slope is positive,
increasing $w_i$ makes things worse, so decrease it. If negative, increase it. That is the
entire idea.
\end{keyidea}

So the update rule writes itself --- move each weight \emph{against} its slope:
\begin{equation}
  \label{eq:gd}
  w_i \;\leftarrow\; w_i \;-\; \eta\,\frac{\partial L}{\partial w_i},
\end{equation}
where $\eta$ (the \textbf{learning rate}) is how big a step to take. This is
\textbf{gradient descent}: walk downhill, one small step at a time.

For our linear $f$ and squared loss, the slope has a form you can compute in your head. With
$f(s) = \sum_i w_i \phi_i(s)$,
\[
  \frac{\partial L}{\partial w_i}
  \;=\; \frac{2}{M}\sum_{m}\bigl(f(s_m) - y_m\bigr)\,\phi_i(s_m) .
\]
Read it in words: \emph{(how wrong I was) $\times$ (how much this feature was present)}. A
feature that was absent gets no blame; a feature that was present in a badly-predicted
position gets a lot.

\begin{worked}{two steps of gradient descent, by hand}
One feature: $\phi_1 =$ material balance in pawns. Three training positions
\emph{(illustrative)}:
\begin{center}
\begin{tabular}{@{}lrr@{}}
\toprule
 & $\phi_1$ & target $y$ \\
\midrule
Position A & $+3$ & $+1.0$ (was won) \\
Position B & $+1$ & $+0.4$ \\
Position C & $-2$ & $-0.8$ \\
\bottomrule
\end{tabular}
\end{center}
Start from $w_1 = 0$ (knowing nothing), learning rate $\eta = 0.05$.

\medskip
\textbf{Step 1.} Predictions $f = 0$ everywhere. Errors $(f - y)$: $-1.0$, $-0.4$, $+0.8$.
\[
\frac{\partial L}{\partial w_1} = \frac{2}{3}\bigl[(-1.0)(3) + (-0.4)(1) + (0.8)(-2)\bigr]
= \frac{2}{3}\bigl[-3 - 0.4 - 1.6\bigr] = \frac{2}{3}(-5.0) = -3.333 .
\]
Negative slope, so increase:
$w_1 \leftarrow 0 - 0.05(-3.333) = +0.1667$.

\medskip
\textbf{Step 2.} Predictions are now $0.500,\ 0.167,\ -0.333$. Errors:
$-0.500,\ -0.233,\ +0.467$.
\[
\frac{\partial L}{\partial w_1} = \frac{2}{3}\bigl[(-0.5)(3) + (-0.233)(1) + (0.467)(-2)\bigr]
= \frac{2}{3}(-2.667) = -1.778 ,
\]
\[
w_1 \leftarrow 0.1667 - 0.05(-1.778) = 0.2556 .
\]

\medskip
The slope shrank from $-3.33$ to $-1.78$: we are approaching the bottom. Two further steps
give $w_1 = 0.30297$ and then $w_1 = 0.32825$, converging towards the exact minimiser
$w_1 = 5/14 = 0.3571$ (Exercise~\ref{ex:ch10:exact} derives it). We have just \emph{learned}
that a pawn is worth about a third of a game result --- from three positions, with no chess
knowledge whatsoever in the algorithm.
\end{worked}

\begin{pitfall}
The learning rate is the whole practical difficulty. Too small and training takes forever; too
large and the step overshoots the valley and lands higher up the opposite side, and the loss
oscillates or explodes. Try $\eta = 1.0$ in the worked example above and watch it diverge
(Exercise~\ref{ex:ch10:lr}). Every "the network diverged" story is this, or something
downstream of it.
\end{pitfall}

\subsection{From a number to a probability}

Chapter~\ref{ch:currency} argued that the output should be a probability. Our $f$ can return
any real number, so we squash it with the logistic function built in §2.3:
\[
  p \;=\; \sigma\bigl(f(s)\bigr) \;=\; \frac{1}{1 + e^{-f(s)}} .
\]
The number $f(s)$ before the squash is called the \textbf{logit}: it lives on the log-odds
scale, where adding is meaningful, and $\sigma$ converts it to a probability.

With a probability output, squared error is the wrong loss --- it is nearly flat when the
model is confidently wrong, which is exactly when you want a strong gradient. The right one
is \textbf{cross-entropy}:
\begin{equation}
  \label{eq:crossent}
  L \;=\; -\bigl[\,y\ln p \;+\; (1-y)\ln(1-p)\,\bigr] .
\end{equation}
If the truth is $y=1$ and you predicted $p = 0.99$, the loss is $-\ln 0.99 = 0.01$. If you
predicted $p = 0.01$, it is $-\ln 0.01 = 4.6$. Confident and wrong is punished savagely,
which is the behaviour you want from anything that will later be asked "how sure are you?"

\begin{keyidea}
Cross-entropy has a beautiful property that will matter in Chapter~\ref{ch:heads}: combined
with the logistic (or softmax) output, its gradient reduces to $(p - y)\phi$ --- literally
\emph{predicted minus actual, times the feature}. The same shape as the linear case. All the
messy calculus cancels.
\end{keyidea}

\section{Where a linear formula hits a wall}

Learning the weights fixed the tuning problem. It did not fix the second failure from
Chapter~\ref{ch:problem}: a sum of terms cannot express \emph{it depends}. Let us prove that
with a chess example precise enough to be provable.

\subsection{A property no linear formula can compute}

Take two binary features of a position:
\[
  \phi_1 = \begin{cases}1 & \text{White castled kingside}\\ 0 & \text{queenside}\end{cases}
  \qquad
  \phi_2 = \begin{cases}1 & \text{Black castled kingside}\\ 0 & \text{queenside}\end{cases}
\]
and ask for a feature every player knows: \textbf{is this an opposite-side-castling
position?} --- the classic marker of a sharp, mutual-attack game.

\begin{center}
\begin{tabular}{@{}ccl@{}}
\toprule
$\phi_1$ & $\phi_2$ & target: opposite sides? \\
\midrule
1 & 1 & 0 (both kingside --- same side) \\
1 & 0 & 1 (\textbf{opposite}) \\
0 & 1 & 1 (\textbf{opposite}) \\
0 & 0 & 0 (both queenside --- same side) \\
\bottomrule
\end{tabular}
\end{center}

\begin{established}
\textbf{No linear formula $w_1\phi_1 + w_2\phi_2 + b$ can produce this pattern.}

\smallskip
\emph{Proof.} Suppose it did, with "produces 1" meaning the output is positive and "produces
0" meaning it is not. The four rows require
\[
\text{(i) } b \le 0,\quad
\text{(ii) } w_1 + b > 0,\quad
\text{(iii) } w_2 + b > 0,\quad
\text{(iv) } w_1 + w_2 + b \le 0 .
\]
Add (ii) and (iii): $w_1 + w_2 + 2b > 0$, so $w_1 + w_2 > -2b$. From (iv),
$w_1 + w_2 \le -b$. Combining: $-b > -2b$, hence $b > 0$ --- contradicting (i). $\blacksquare$
\end{established}

This is the classic XOR argument, and it is not a curiosity: \textbf{"it depends" is XOR}.
Every piece of chess knowledge of the form \emph{this is good, unless that, in which case it
is bad} has this shape, and a weighted sum of features is structurally incapable of holding
it. You cannot fix it by tuning harder. You cannot fix it with more data. The form of the
equation is wrong.

Classical engines coped by hand-coding the interaction as a new feature --- literally adding
$\phi_3 = $ "opposite-side castling" with its own weight. That works exactly once per
interaction somebody thinks of, and there are more interactions in chess than there are
chess programmers.

\subsection{The repair: one hidden layer}

Instead of going straight from features to output, go through an intermediate stage of
\emph{learned} quantities, with a non-linearity in between. The non-linearity of choice is
the simplest one imaginable:
\begin{equation}
  \label{eq:relu}
  \relu(z) \;=\; \max(0, z) \qquad \text{"negative becomes zero, positive passes through".}
\end{equation}

Now build the thing that was just proved impossible.

\begin{worked}{a two-neuron network that computes ``opposite-side castling''}
Hidden units:
\[
  h_1 = \relu(\phi_1 + \phi_2 - 0), \qquad h_2 = \relu(\phi_1 + \phi_2 - 1),
\]
output:
\[
  f = h_1 - 2h_2 .
\]
Evaluate all four cases:
\begin{center}
\begin{tabular}{@{}ccrrrl@{}}
\toprule
$\phi_1$ & $\phi_2$ & $h_1$ & $h_2$ & $f = h_1 - 2h_2$ & target \\
\midrule
0 & 0 & $\relu(0)=0$ & $\relu(-1)=0$ & $0$ & 0 \checkmark \\
1 & 0 & $\relu(1)=1$ & $\relu(0)=0$  & $1$ & 1 \checkmark \\
0 & 1 & $\relu(1)=1$ & $\relu(0)=0$  & $1$ & 1 \checkmark \\
1 & 1 & $\relu(2)=2$ & $\relu(1)=1$  & $2 - 2 = 0$ & 0 \checkmark \\
\bottomrule
\end{tabular}
\end{center}
Four for four. Two neurons and a $\max(0,\cdot)$ have expressed a piece of chess knowledge
that no weighted sum of the same two features can express.
\end{worked}

Look at what the hidden units became. $h_1$ fires when \emph{at least one} king went kingside;
$h_2$ fires only when \emph{both} did. Neither was asked for. They are \textbf{learned
features} --- intermediate concepts invented because they were useful for the final answer.
That is the whole idea of a hidden layer, and everything Part~\textsc{v} tries to read out of
Leela's middle layers is a much larger version of exactly this.

\begin{keyidea}
\textbf{A network is a stack of the Chapter~\ref{ch:problem} formula.} Each layer computes a
weighted sum of the layer below and squashes it:
\[
  \vv{x}^{(l)} = \relu\bigl(W^{(l)}\vv{x}^{(l-1)} + \vv{b}^{(l)}\bigr).
\]
The first layer's inputs are the raw board; the last layer's output is the evaluation; and
every layer in between is a set of features that the training process found worth having. The
$\relu$ is what makes stacking worthwhile --- without it, a stack of linear maps collapses
into a single linear map (Exercise~\ref{ex:ch10:collapse}), and we are back where we started.
\end{keyidea}

\section{A forward pass, by hand}

Here is a complete miniature network, so the words above have arithmetic under them.

\begin{worked}{one forward pass}
Two inputs, two hidden units, one output. \emph{(Illustrative weights.)}
\[
W^{(1)} = \begin{pmatrix} 0.5 & -1.0 \\ 2.0 & 0.5\end{pmatrix},\quad
\vv{b}^{(1)} = \begin{pmatrix} 0.1 \\ -0.4 \end{pmatrix},\quad
W^{(2)} = \begin{pmatrix} 1.5 & -0.8\end{pmatrix},\quad b^{(2)} = 0.2 .
\]
Input $\vv{x} = (1.0,\ 0.6)$.

\smallskip
\textbf{Hidden pre-activations:}
\[
z_1 = 0.5(1.0) + (-1.0)(0.6) + 0.1 = 0.5 - 0.6 + 0.1 = 0.0,
\]
\[
z_2 = 2.0(1.0) + 0.5(0.6) - 0.4 = 2.0 + 0.3 - 0.4 = 1.9 .
\]
\textbf{After ReLU:} $h_1 = \max(0, 0.0) = 0.0$, $h_2 = \max(0, 1.9) = 1.9$.

\smallskip
\textbf{Output logit:} $f = 1.5(0.0) + (-0.8)(1.9) + 0.2 = -1.52 + 0.2 = -1.32$.

\smallskip
\textbf{As a probability:} $\sigma(-1.32) = \dfrac{1}{1+e^{1.32}} = \dfrac{1}{1+3.743} = 0.211$.

\smallskip
So this network says: expected score $0.211$, i.e.\ $V = 2(0.211) - 1 = -0.578$. Note that
$h_1$ contributed nothing --- its input landed exactly at the ReLU's kink. That is normal:
in a large network most units are silent for any given input, and \emph{which} units are
silent is itself information about the position. Chapter~\ref{ch:probes} is about reading it.
\end{worked}

\section{Training, in one page}

The remaining machinery, stated rather than derived, because you will not need to implement
it:

\begin{enumerate}
\item \textbf{Backpropagation} is gradient descent applied to a stack. To know how a weight in
  layer 3 affects the loss, you multiply together the effect of layer 3 on layer 4, layer 4 on
  layer 5, and so on --- the chain rule. Doing this efficiently, once, for all weights, is
  backpropagation. It is bookkeeping, not a new idea.
\item \textbf{Batches.} Computing the gradient over millions of positions per step is
  wasteful; use a random few hundred each step. The gradient is then noisy, which turns out to
  help.
\item \textbf{Initialisation} matters: all-zero weights make every unit identical forever, so
  weights start random and small.
\item \textbf{The loss usually has several terms.} Leela's has three (Chapter~\ref{ch:heads}):
  one for the value head, one for the policy head, one for the moves-left head, plus a
  regularisation term that discourages large weights.
\end{enumerate}

\begin{established}
Two things worth knowing about what training does and does not guarantee. It finds a
\emph{local} minimum, not the best possible network --- in practice, for large networks, this
matters far less than intuition suggests. And it optimises the loss on the training data,
which is only useful insofar as that data resembles what the network will meet later. Leela's
training data is its own games, which is why its policy models \emph{Leela's} choices rather
than a human's --- a fact with direct consequences for how you read a policy prior in a
training tool for humans.
\end{established}

\begin{repobox}[title={in this repository}]
You never train the main network here --- \texttt{engine/*.pb.gz} are downloaded, trained by
the Leela project on thousands of GPUs. But you \emph{will} train small models: a linear
probe is exactly the model of §10.2 (Chapter~\ref{ch:probes}), and the salience ranker that
learns which board facts a master would mention is a slightly larger one
(Chapter~\ref{ch:speak}). Everything in this chapter applies to them directly.
\end{repobox}

\begin{exercises}

\exhead[one more step]{ch10:onestep} Continue the worked example of §10.2.2 for a third and
fourth step of gradient descent. Does the slope keep shrinking? Sketch $w_1$ against step
number.

\exhead[the exact answer]{ch10:exact} For a one-feature linear model with squared loss, the
minimiser is $w_1 = \sum_m \phi_m y_m / \sum_m \phi_m^2$. Compute it for the three positions
in the example and check the claimed $5/14$. How many gradient steps of §10.2.2 are needed to
get within $0.05$ of it? Why does gradient descent approach the answer ever more slowly?

\exhead[learning rate]{ch10:lr} Redo two steps of the same example with $\eta = 1.0$. What
happens to $w_1$? Now try $\eta = 0.3$. Describe the three regimes (too small, about right,
too large) in one sentence each.

\exhead[blame assignment]{ch10:blame} The gradient is (error) $\times$ (feature). Explain why
a feature that is zero in a position receives no update from it, and why that is the right
behaviour. What does this imply for a feature that is almost always zero?

\exhead[cross-entropy]{ch10:crossent} Compute the cross-entropy loss for $y = 1$ with
$p = 0.9, 0.5, 0.1, 0.01$. Plot roughly. Compare with the squared errors for the same four.
Which loss cares more about confident mistakes, and by how much at $p = 0.01$?

\exhead[the XOR proof]{ch10:xorproof} Reproduce the impossibility proof of §10.3.1 without
looking. Then find another chess property with the same logical shape (good, \emph{unless}),
and state it as a truth table over two binary features.

\exhead[build the network]{ch10:buildxor} Verify the two-neuron network of §10.3.2 on all four
inputs yourself. Then modify it so it outputs $+1$ for same-side castling and $0$ for
opposite --- changing as few numbers as possible.

\exhead[why non-linearity]{ch10:collapse} Show that two linear layers without a ReLU collapse
to one: if $\vv{y} = W_2(W_1\vv{x} + \vv{b}_1) + \vv{b}_2$, find the single $W$ and $\vv{b}$
with $\vv{y} = W\vv{x} + \vv{b}$. What does this say about the value of depth without
non-linearity?

\exhead[forward pass]{ch10:forward} Run the network of §10.4 on the input $\vv{x} = (0.2,
-1.0)$. Give $z_1, z_2, h_1, h_2$, the logit, the probability and $V$. Which hidden unit is
silent this time?

\exhead[counting parameters]{ch10:params} A layer mapping 768 numbers to 768 numbers has how
many weights (ignore biases)? Leela's BT3 has 15 such layers plus attention machinery.
Estimate the parameter count to within a factor of two, and compare with the three weights
you fitted by hand in §10.2.2.

\exhead[what the target really is]{ch10:target} Leela's value target is the result of its own
self-play game, not a human game. Give one way this makes the network stronger and one way it
makes its policy a slightly wrong model of \emph{your} chess --- and say which chapter's
metric is affected by the latter.

\end{exercises}
